% Part: first-order-logic
% Chapter: axiomatic-deduction
% Section: proving-things-quant

\documentclass[../../../include/open-logic-section]{subfiles}

\begin{document}

\olfileid{fol}{axd}{prq}

\olsection{संख्यापकांसह \usetoken{P}{derivation}}

\begin{ex}
$(\lforall[x][!A(x)] \land
\lforall[y][!B(y)]) \lif \lforall[x][(!A(x) \land !B(x))]$ ची
!!a{derivation} देऊ.

प्रथम, लक्षात घ्या की
\begin{align*}
  (\lforall[x][!A(x)] \land \lforall[y][!B(y)]) & \lif \lforall[x][!A(x)]\\
  \intertext{हे \olref[prp]{ax:land1} चे उदाहरण आहे, आणि}
  \lforall[x][!A(x)] & \lif !A(a) \\
  \intertext{हे \olref[qua]{ax:q1} चे उदाहरण आहे. म्हणून \olref[pro]{prop:chain} वरून आपल्याला माहीत आहे की}
  (\lforall[x][!A(x)] \land \lforall[y][!B(y)]) & \lif !A(a) \\
  \intertext{हे !!{derivable} आहे. त्याचप्रमाणे,}
  (\lforall[x][!A(x)] \land \lforall[y][!B(y)]) & \lif \lforall[y][!B(y)] \qquad\text{आणि}\\
    \lforall[y][!B(y)] & \lif !B(a)\\
    \intertext{ही अनुक्रमे \olref[prp]{ax:land2} आणि \olref[qua]{ax:q1} यांची उदाहरणे असल्यामुळे,}
    (\lforall[x][!A(x)] \land \lforall[y][!B(y)]) & \lif !B(a)\\
    \intertext{हे \olref[pro]{prop:chain} मुळे सिद्धतायोग्य आहे. \olref[prp]{ax:land3} चे योग्य उदाहरण आणि~\MP चे दोन उपयोग करून आपल्याला दिसते की}
    (\lforall[x][!A(x)] \land \lforall[y][!B(y)]) & \lif (!A(a) \land !B(a))\\
    \intertext{हे सिद्धतायोग्य आहे. आता \QR{} लावून पुढील सूत्र मिळते:}
(\lforall[x][!A(x)] \land \lforall[y][!B(y)]) & \lif \lforall[x][(!A(x) \land !B(x))].
\end{align*}
\end{ex}

\end{document}
